\chapter{Gastritis Medicines}
\index{Gastritis Medicines}

\section*{Definition and Overview}
Medicines are central to the treatment of gastritis, the inflammation of the stomach lining. They work in three main ways: reducing the acid that damages the lining, eradicating the H.\ pylori bacterium that causes chronic gastritis, and protecting or healing the stomach lining itself. The main medicines are proton pump inhibitors (PPIs) such as omeprazole and esomeprazole, H2 blockers such as ranitidine, antacids, and the antibiotics used in H.\ pylori eradication. Choosing the right medicine depends on the cause of the gastritis and its severity. This chapter explains the medicines used for gastritis, how they work, and how they are taken.

\section*{Epidemiology}
Acid-suppressing medicines are among the most prescribed drugs in many countries and worldwide. Proton pump inhibitors are taken by around 10--15% of people, often long term, and they are the mainstay of gastritis and ulcer treatment. H.\ pylori, which causes chronic gastritis in around a third of people, is treated with a 7--14 day course of antibiotics combined with a PPI, which eradicates the infection in around 85% of people. While these medicines are very effective, long-term PPI use is associated with small risks, so they are used at the lowest effective dose and reviewed regularly.

\section*{Aetiology and Pathophysiology}
\begin{itemize}
    \item \textbf{Acid in gastritis:} inflammation damages the stomach lining, and acid worsens it, so reducing acid allows healing
    \item \textbf{Proton pump inhibitors (PPIs):} block the enzyme in stomach cells that produces acid, providing the strongest acid suppression; used for 4--8 weeks for healing gastritis and ulcers
    \item \textbf{H2 receptor blockers:} reduce acid by blocking histamine signals to acid-producing cells; slightly weaker than PPIs and used for milder disease
    \item \textbf{Antacids:} neutralise acid already in the stomach, giving rapid but short-lived relief
    \item \textbf{H.\ pylori eradication:} antibiotics (often two, such as amoxicillin and clarithromycin) combined with a PPI for 7--14 days; the PPI improves antibiotic effectiveness
    \item \textbf{Cytoprotective medicines:} such as bismuth and sucralfate, which coat and protect the lining, used in some regimens
    \item \textbf{Stopping the cause:} withdrawing NSAIDs or alcohol is as important as any medicine
    \item \textbf{Long-term use:} PPIs are continued long term in some people at risk of ulcers, such as those who must take NSAIDs
\end{itemize}

\section*{Clinical Features}
The relevant features are how the medicines are used.
\begin{itemize}
    \item PPIs taken once daily, usually before breakfast, for the prescribed course
    \item Improvement in pain and nausea within days to weeks
    \item H.\ pylori treatment taken twice daily for 7--14 days, with the full course completed
    \item Antacids used as needed for rapid symptom relief
    \item Common side effects: headache, nausea, and bowel changes
    \item The need for review, since long-term PPI use should be reassessed
    \item Re-testing for H.\ pylori after treatment to confirm eradication
    \item The return of symptoms if medicines are stopped too early or the cause is not removed
\end{itemize}

\section*{Differential Diagnosis}
Medicine selection depends on the diagnosis.
\begin{itemize}
    \item H.\ pylori-positive gastritis, which needs antibiotics plus a PPI
    \item NSAID-related gastritis, which needs a PPI and stopping or changing the NSAID
    \item Reflux symptoms, which may respond to PPIs or H2 blockers
    \item Bleeding gastritis or ulcers, which need urgent treatment with high-dose PPIs
    \item Functional dyspepsia, where medicines help less and lifestyle matters more
    \item The risk of long-term PPI use, including a small increased risk of fractures, infections, and low magnesium
\end{itemize}

\section*{When to Seek Medical Care}
Anyone with symptoms of gastritis should see a doctor before self-medicating long term. Medical review is needed for persistent symptoms, black stools, vomiting blood, or the need for long-term acid suppression.

\textbf{Contact your local emergency services for:}
\begin{itemize}
    \item Vomiting blood or black, tarry stools
    \item Severe abdominal pain
    \item Faintness, collapse, or signs of significant bleeding
    \item Severe allergic reactions to medicines
    \item Chest pain or breathlessness
\end{itemize}

\section*{Diagnosis and Investigations}
\begin{itemize}
    \item \textbf{H.\ pylori testing:} before treatment, using a breath, stool, or biopsy test
    \item \textbf{Endoscopy:} for persistent symptoms, bleeding, or warning features, with biopsy
    \item \textbf{Blood tests:} for anaemia and, in long-term PPI users, periodic monitoring where indicated
    \item \textbf{Medicine review:} identifying NSAIDs and other drugs that cause or worsen gastritis
    \item \textbf{Post-treatment testing:} confirming H.\ pylori eradication at least 4 weeks after finishing antibiotics
\end{itemize}

\section*{Management}
\begin{itemize}
    \item \textbf{PPIs:} omeprazole, esomeprazole, pantoprazole, or rabeprazole, once daily before food for 4--8 weeks
    \item \textbf{H2 blockers:} for milder disease or when PPIs are not suitable
    \item \textbf{Antacids:} for rapid relief of mild symptoms, taken as needed
    \item \textbf{H.\ pylori triple therapy:} a PPI with two antibiotics for 7--14 days, repeated with a different combination if first-line treatment fails
    \item \textbf{Bismuth quadruple therapy:} for resistant infection, including bismuth, a PPI, and two antibiotics
    \item \textbf{Prevention during NSAID use:} a low-dose PPI for people who must continue NSAIDs
    \item \textbf{Stop the cause:} reduce alcohol, stop smoking, and review NSAIDs
    \item \textbf{Long-term review:} the lowest effective dose of PPI, with regular reassessment of the need
\end{itemize}

\section*{Complications}
\begin{itemize}
    \item Side effects of medicines, including headache, diarrhoea, and nausea
    \item Clostridioides difficile infection with antibiotics, rarely
    \item Long-term PPI risks: small increases in fractures, intestinal infections, and low magnesium
    \item Interactions, including PPIs with some blood thinners and other drugs
    \item Eradication failure, requiring a second course
    \item Progression of untreated or resistant gastritis to ulcers and bleeding
\end{itemize}

\section*{Prognosis and Outcomes}
Gastritis medicines are highly effective. Symptoms settle within weeks of acid suppression, H.\ pylori is eradicated in around 85% of people, and ulcers and bleeding are prevented and treated. Most people complete a short course and stop, while those at risk continue the lowest effective dose. The key to the best outcomes is accurate diagnosis, treatment of H.\ pylori when present, removal of the cause, and review of long-term medicine use.

\section*{Prevention and Self-Care}
\begin{itemize}
    \item Take medicines as prescribed and complete the full course of antibiotics
    \item Take PPIs before meals as directed for best effect
    \item Avoid NSAIDs and alcohol while the stomach heals
    \item Confirm H.\ pylori eradication with a follow-up test
    \item Attend review appointments to reassess the need for long-term medicine
    \item Report black stools, vomiting blood, or persistent symptoms
    \item Do not self-prescribe acid suppressors for months without medical review
\end{itemize}

\section*{Sources and Further Reading}
The following reputable sources provide further information for healthcare students and patients seeking reliable, current advice about gastritis medicines.
\begin{itemize}
    \item Healthdirect Australia. \textit{Gastritis treatment.} https://www.healthdirect.gov.au/
    \item The Gut Foundation. \textit{Gastritis and medicines.} https://www.gutfoundation.org.au/
    \item Therapeutic Goods Administration. \textit{Medicine information.} https://www.tga.gov.au/
    \item NHS. \textit{Gastritis treatment.} https://www.nhs.uk/conditions/gastritis/
    \item Mayo Clinic. \textit{Gastritis treatment.} https://www.mayoclinic.org/
    \item MedlinePlus. \textit{Gastritis.} https://medlineplus.gov/
\end{itemize}
